\newcommand{\wineIntro}{Tonight, to fill us with extra inspiration and zeal, we take a deeper look at just how tantalizingly close we are today to finally achieving these generational goals of bringing freedom to everyone.}

\newcommand{\wineSummaryThisYear}{This year, we will celebrate just how close worldwide freedom really is, including information from the book \textit{Abundance} by Peter Diamandis and Steven Kotler.}

\newcommand{\wineFirstCup}{
  \newReading With this first cup of wine, we consider just how unique the present is in the scope of history, and why that means we may finally achieve the dream of bringing freedom to the whole human race.
  
  From what we see in the news every day, it may seem like the world is going to pieces around us, and that fighting for good causes and progress is like a game of whack-a-mole. But what if that is just a misconception; one that is easy to believe because of the way our brains are wired.
  
  Imagine a world of nine billion people with clean water, nutritious food, affordable housing, personalized education, top-tier medical care, and nonpolluting, ubiquitous energy. 

  There are serious people who believe that those goals are all achievable by 2035, with noticeable change happening in the next decade.

  The root of what has these people convinced is that sometime during the last 150 years, the rate of technological progress began doing something unprecedented in all of human history up to that point; we started experiencing exponential growth.
  
  \newReading Over the past 100,000 years, Homo sapiens evolved in a ``local and linear'' world. In our ancestor's local environment, most everything that occurred in their day happened within a day's walk. In their linear environment, change was excruciatingly slow\textemdash{}life from one generation to the next was effectively the same\textemdash{}and what change did arrive always followed a linear progression. To give you a sense of the difference, if I take thirty linear steps from my front door, I end up 100 feet down the street. But, if I take thirty exponential steps (one, two, four, eight...) I end up three billion feet away, effectively lapping the globe twenty-six times.  Exponential growth is how we reached the point in 2013 where we were producing more digital information every ten minutes than had been created from the beginning of time up until 2003.

  The global trends and exponential growth that characterize today's world are incredibly difficult to comprehend because we've simply never seen it before. Five hundred years ago, technologies were not doubling in power and halving in price every eighteen months. Waterwheels were not becoming cheaper every year. A hammer was not easier to use from one decade to the next. The yield of corn seed varied by the season's climate, instead of improving each year. Every twelve months, you could not upgrade your oxen's yoke to anything much better than what you already had.
  
  Exponential growth has taken hold in a variety of different fields, but perhaps the one we are most familiar with is the field of computing power. A classic example is that our smartphones are thousands of times more powerful than the computers NASA used to get to the moon. While we can see looking backwards that was an incredible rate of growth, it is still challenging to realize what exponential growth in computing means the future will bring. It is hard for our brains to grasp that we are still on track, from predictions made decades ago, for a \$1,000 laptop to have the same computing power as a human brain by 2025.
  
  \newReading For this Seder, we will explore the impact exponential growth will have on bringing unprecedented freedom to the world, and ways you can help accelerate it even further. Realizing that we've entered a point in history where worldwide freedom and abundance are actually achievable near-term goals, and the future is not a bleak landscape of scarcity and despair, may even inspire us to focus on and accomplish these changes even sooner.
  
  Thinking the world is getting worse and disasters are imminent is common. There are strong cognitive biases built into our brains that lead to those thoughts. And those thoughts are not a new trend that emerged in modern society; in 1828 John Stuart Mill said, ``I have observed that not the man who hopes when others despair, but the man who despairs when others hope, is admired by a large class of persons as a sage.''
  
  With this first cup of wine, we attempt to put our inherent biases aside during this Seder and listen openly to the thoughts presented from people who hope.
}

\newcommand{\wineSecondCup}{
  \newReading With this second cup of wine, we look at how close we are to a world where everyone is free to live out a full life.
  
  The next fifteen years will not result in a world where everyone lives in a penthouse apartment and drives a Tesla. But what is possible is providing everyone on this planet with a life of possibility; a life where the basics are covered, and then some. A world where everyone's days are spent dreaming and doing, not struggling to merely survive.
  
  \subsubsection*{The Scarcity Mindset}
  Many people's immediate reaction to such a bold claim is that there simply aren't enough resources in the world to achieve these goals, but resource scarcity is most often a matter of perspective.
  
  Today, aluminum is one of the cheapest things we can buy, and is regularly used once and then thrown away. But as recently as the 1800s, it was literally the most valuable metal in the world; Napoleon III threw a banquet for the king of Siam where the honored guests were given aluminum utensils, while the others had to make do with gold. Chemically speaking, aluminum is incredibly abundant, making up 8\% of the weight of the Earth's crust, but it wasn't until 1886 that a technology was invented to cheaply purify aluminum from the clay in which it is commonly found.
  
  \newReading The story of aluminum is not unusual; history is littered with tales of once-rare resources made plentiful by innovation. Imagine an orange tree packed with fruit. If I pluck all the oranges from the lower branches, I am effectively out of accessible fruit. From my limited perspective, oranges are now scarce. But once someone invents a piece of technology called a ladder, I've suddenly got new reach. Problem solved. Technology is a resource-liberating mechanism. It can make the once scarce now abundant. 
  
  A recent example is the dramatic drop in US gasoline prices from 2012\textendash{}2016 due to innovations in extraction technology that made oil in tar sands and other previously inaccessible areas suddenly available. Or take India in 2002, which could not produce enough cotton to meet its own needs and had to import it. But after the introduction of the engineered strain Bt-cotton, which grew abundantly, India became a net exporter of cotton by 2011. Lack of access to basic health care and electricity are two areas restricting people's freedom to live a full life that are not truly cases of scarce resources, but merely inaccessible ones.  
  
  \subsubsection*{Basic Health}
  \newReading In 2000, two people were dying almost every minute from malaria\textemdash{}nearly one million each year. The medicines to combat this were not a scarce resource (twelve-cent drugs could prevent half of all malarial childhood deaths), but lack of diagnostics and distribution made them inaccessible. Bill Gates has taken his fortune, generated due to the enormous exponential technological growth in computing, and applied it to medical and other technologies to accelerate their growth. The results have been another phenomenal example of accelerating exponential growth. From 2000\textendash{}2006, the number of deaths dropped by around 2\% each year. From 2006\textendash{}2009 the number dropped by 4\% each year. From 2009\textendash{}2016 the number dropped by 8\% each year. For the latest numbers from 2018, 405,000 people died from malaria. That is still far too high, but it is well under half the number from 18 years prior.
  
  Similar trends have been happening with other areas of health. In 1990, 35,000 children under the age of five were dying every day; out of every 1000 live births, 90 of those babies would never celebrate their fifth birthday. By 2016, the mortality rate was cut in half\textemdash{}20,000 kids every day have been given the freedom to keep on living instead of having their lives cut tragically short.
  
  \subsubsection*{Clean Energy}
  \newReading A lack of energy is one of the strongest chains holding people back from living a free life. 3.5 billion people rely on primitive fuels such as wood or charcoal for cooking and heating, creating an enormous ripple effect. Women and children who spend long periods of time around traditional open fires inhale the equivalent of two packs of cigarettes a day, leading to pneumonia, bronchitis, lung cancer, and death. Because children have to spend daylight hours helping search for fuel, time spent on their education is severely reduced, and it's very challenging to do homework or study at night without lighting. And because many teachers don't want to work in communities without lights or electricity, schools are understaffed and less effective. But once again, energy is not truly scarce, it is merely inaccessible.  Five thousand times more solar energy reaches the Earth's surface than all of humanity currently uses. Fortunately, just like computing power, since 1976 solar technology has been on an exponential growth curve, halving in price every ten years.
  
  One aspect that causes our brains to fail to grasp the implications of exponential growth is that the effects start out moving seemingly very slowly. Who cares if from 1976\textendash{}1985 the cost of solar dropped from 100 times the cost of coal down to fifty times the cost\textemdash{}it was still incredibly expensive and impractical. But the amazing thing about exponential growth is that it \textit{keeps} accelerating. In 2009, solar power was still roughly 2.5 times the cost of fossil fuel energy.  But the price keeps dropping; between 2014\textendash{}2018, country after country and state after state have passed the point of ``grid parity''\textemdash{}where the cost of solar electricity is less than the cost of electricity from fossil fuels. The battery technology needed to store power for when the sun isn't shining is also on an exponential trend.
  
  \newReading When critics point out that solar still only accounts for 2\% of the world's total electrical production, that's linear thinking in an exponential world. Solar panel installation has been growing by more than 30\% annually since the early 1990s; even if it slows down a little bit, that means we're less than fifteen years from having enough capacity to supply all of the world's electrical needs. With abundant and cheap energy, you can skyrocket economic growth (the US spends more than \$1 trillion on energy each year at current prices), supply everyone on the planet with clean purified water, scrub CO\textsubscript{2} out of the atmosphere, and refrigerate vaccines and food throughout the developing world, just to name a few possibilities.
  
  We drink this second cup with a better understanding of the incredible improvements in technology that will give so many more people the freedom to live. 

}

\newcommand{\wineThirdCup}{
  \newReading With the third cup of wine, we learn of advances that will give people not just the freedom to survive, but the freedom to thrive.
  
  \subsubsection*{Accessible Education}
  Education has long been recognized as instrumental in increasing the amount of freedom a person has to thrive; it is the key to health and nutrition, to overall improvements in standard of living, to better agricultural and environmental practices, and to higher gross national product.
  
  Teachers are in short supply, especially in the developing world, but exponential technological growth is giving rise to the possibility of self-directed, effective learning independent of formal instruction. The first step in this process was the tremendous rise in the literacy race. From 1800 to 1900 it improved slightly from 12\% up to 22\%, but then exponential growth kicked in and by 1960 half the world could read and today the literacy rate is up to 86\%. Now that literacy is becoming nearly universal, the opportunities for further education are opening wide.
  
  \newReading From 1999\textendash{}2005, Sugata Mitra conducted a series of experiments in rural areas and city slums in various countries demonstrating that groups of children, irrespective of who or where they are, can learn to use computers and the Internet on their own with public computers in open spaces such as roads and playgrounds, even without knowing English. In two months, impoverished twelve-year-olds in southern India used a computer to teach themselves enough about genetics and biotechnology to score half as well on an exam as high school kids studying biotech at the best schools in New Delhi. As a follow-up, a slightly older girl\textemdash{}who knew nothing about biotechnology\textemdash{}was recruited to ``tutor'' them using the grandmother method: just stand behind the kids and provide encouragement, ``Wow, that's fantastic, show me something else!'' Incredibly, after two more months, the twelve-year-olds were scoring just as well on tests as the New Delhi high schoolers.
  
  With plummeting prices of computers and expanding internet coverage, installation of these portals to facilitate ``self-organized learning environments'' can soon begin in earnest. And we've already gotten a four-fold head start\textemdash{}forget One Laptop Per Child, people have found that in these environments ``one child in front of a computer learns little; four discussing and debating learn a lot.''
  
  The advent and growth of online video tutorials is also accelerating this process. The Khan academy had 50,000 people watching their videos each month in 2009, in 2011 it was two million per month, in 2017 it was 40 million per month. These learners are spread across 200 different countries, watching videos in 36 different languages, all for free. Exponential growth hard at work yet again.
  
  \subsubsection*{The End of Poverty}
  \newReading In 1820, 94\% of the world lived on the equivalent of less than two dollars per day (in today's money). That's what the WHO defines as ``extreme poverty'' and certainly makes it hard to live a full and free life. In 1920, that was down to 88\% but then exponential change kicked in: 60\% in 1970, 30\% in 2000, and all the way down to 10\% in 2017.
  
  What has driven this change has been increasing access to technology. In 2016, there were an astonishing 82 cell phones for every 100 people in Africa. Mobile banking allows the global poor to have access to financial services and transactions that most of us take for granted: safely receiving a paycheck and storing it for later, applying for micro-loans, sending money to friends and family in need. Cell phones allow people to find out where people are who want to buy what they are selling; a famous example of cell phones penetrating a fishing community in southern India\textemdash{}allowing fishermen to know which port still needed fish and which were already full for the day\textemdash{}resulted in a 6\% increase in the fishermen's profits and a 4\% drop in the cost to consumers. Studies estimate that just cell phones alone can add 0.5\% to a country's GDP growth rate for each new 10\% of the population that gets access to them.
  
  \subsubsection*{Freeom to Pursue Happiness}
  \newReading Studies have found that money can buy happiness\textemdash{}up to a point. Americans' emotional satisfaction rises in lock-step with income until it reaches \$75,000 per year. Surprisingly, at that point the correlation disappears. The freedom to flourish\textemdash{}to truly enjoy a life of possibility\textemdash{}costs roughly \$75,000 per year...in the United States.  For the average global citizen, that amount is only \$10,000. Now ten grand is certainly not chump change, but exponential growth is driving that number downwards through the forces of demonetization and dematerialization.
  
  Dematerialization is a radical decrease in the size and cost for items we use in our lives. Twenty years ago, it used to be that a ``happy'' American owned a camera, a CD player, a stereo, a watch, an alarm clock, a set of encyclopedias, and a whole bunch of other assets, which could easily add up to more than \$10,000. Poof! All of those items have disappeared today and instead exist inside your cell phone at a tiny fraction of the cost. If you take even just a few of the features that come standard on a \$100 smart phone today, and add up the prices that those technologies cost when they first came out in the 1980s (video conferencing, GPS, voice recording, video players, gaming consoles, a digital watch...), you easily surpass \$1 million worth of technology that has been dematerialized. We are getting the same benefit to our happiness, but at a fraction of the cost.
  
  \newReading Demonetization is when technology makes a service or item so inexpensive that it is virtually or literally free; money is no longer needed for someone to enjoy the associated happiness. Craigslist demonetized classified ads, taking 99\% of the profits out of the newspaper industry and putting them back into the pockets of the consumer. Or eBay, which demonetized transactions, increasing the availability of goods while simultaneously reducing their cost. In 2019 there were 2.8 million apps available for free in the Google Play store.
  
  With this third cup of wine, we realize that with access to more education, healthcare, and energy raising everyone's incomes, while dematerialization and demonetization are lowering the price tag for achieving happiness, it will not be long before these two lines intersect. The freedom for everyone in the world to have a thriving life of potential and possibility is near at hand.
}

\newcommand{\wineFourthCup}{
  \newReading As our Seder draws near to the end, we take up our cups one last time as we think about how we can further accelerate the exponential growth to freedom.
  
  \subsubsection*{Driving Innovation}
  While universal freedom is a very real possibility, we're also in a race against time. Until the innovations of abundance bear fruit, scarcity remains a real concern. And nearly as bad as scarcity is the threat of scarcity and the devastating violence it can often incite. How can we steer innovation and step on the gas?
  
  Decades ago, Buckminster Fuller observed that ``You never change things by fighting the existing reality. To change something, build a new model that makes the existing model obsolete.'' In our current era of exponential technological growth, where we are so close on so many fronts to achieving universal freedom, many of the most high-profile philanthropists have taken Fuller's quotation to heart and changed how they donate their money. Rather than solely\textemdash{}and sometimes even instead of\textemdash{}supporting traditional charities or NGOs, they are increasingly supporting organizations that spur radical innovation aimed at addressing societal problems.
  
  \newReading The most famous of these organizations is the X PRIZE Foundation, which launches challenges offering large cash prizes to any team who can solve a pressing problem; automobile fuel efficiency, global learning, women's safety, and atmospheric water extraction are just some of the goals they have set. One of the most successful of these challenges was launched in the wake of the \textit{Deepwater Horizon} BP oil spill. Shocked at the magnitude of environmental damage that occurred, philanthropist Wendy Schmidt agreed to bankroll the competition. The focus of the prize was clear: the technology used to clean up the BP spill in 2010 was the same as was used to clean up the Exxon Valdez spill in 1989\textemdash{}it was clearly time for an upgrade. Less than 18 months after the BP oil spill, a team won the challenge with a new process that had quadrupled the performance of the oil industry's existing cleanup technology. What made this even more exciting was that this giant breakthrough had been achieved just by offering a \$1 million prize. Considering that the environmental damage of just that single spill was estimated at \$17 billion, the ability to prevent environmental damage by cleaning up all future spills more effectively made the return on charitable investment enormous.
  
  People across the spectrum are getting behind the concept of incentive prizes as an extremely effective method of achieving social change. Software entrepreneurs such as Elon Musk, Larry Page, and Sergey Brin, motivational speaker Tony Robbins, film director James Cameron, columnist Arianna Huffington, and former first lady Barbara Bush are just a few of the strong supporters of the X PRIZE.
  
  \newReading Frustrated after an oil company managed to secure drilling rights to land in the Amazon Rainforest\textemdash{}that a charity he had donated to was supposed to protect\textemdash{}philanthropist Vishen Lakhiani switched his support to the X PRIZE Foundation and explained his reasoning in a company memo: ``Why keep fighting companies acre by acre through the rainforest when we can work to make oil unnecessary with solar power innovation? Or make food production so efficient that rainforest doesn't need to be cut down.''

  You don't have to be rich enough to sponsor your own X PRIZE to make a difference. Anyone can decide to donate to support part of the next great humanitarian breakthrough.
  
  \subsubsection*{Don't Fall Victim to Cognitive Biases}
  The more people who understand the unprecedented opportunity for freedom that the near future holds, the faster we'll achieve this age-old goal. It can be tempting to dismiss all these thoughts and ideas about the potential for universal freedom and abundance as na\"ive optimism; our brains are wired to jump to that conclusion. But don't let your biology control you. Recognize these biases that are luring you to reject this exciting future, and point them out to others who fall into their trap.
  
  Confirmation bias is a tendency to search for or interpret information in a way that confirms one's preconceptions\textemdash{}but it can often limit our ability to take in new data and change old opinions. This bias is typified in the huge uptick in downright lies told by politicians that still persist in large portions of the public's minds because people are able to find so many other ``sources'' on the internet or cable news networks that repeat that false information. If you are thinking that there's no way we can free the poor because the hole we're in is too deep to climb out of, check that you are actually giving new information the same weight as what you already ``know.''
  
  \newReading Anchoring bias is the predilection for relying too heavily on one piece of information when making decisions. When people believe the world's falling apart, it's often an anchoring problem. At the end of the 19\textsuperscript{th} century, London was becoming uninhabitable because of the accumulation of horse manure. People were absolutely panicked. Because of anchoring, they couldn't imagine any other possible solutions. No one had any idea the car was coming and soon they'd be worrying about dirty skies, not dirty streets.
  
  Just because current water purification technology requires too much electricity to deploy effectively in developing countries shouldn't make you forget that abundant energy is coming in the next few years that will completely alter those costs.
  
  The negativity bias\textemdash{}the tendency to give more weight to negative information and experiences than positive ones\textemdash{}sure isn't helping matters either. That bias gets further compounded by the strategy of news sources to follow the ``if it bleeds, it leads'' philosophy of almost exclusively reporting violent and negative news because it gets higher ratings; often over 90\% of articles are pessimistic.
  
  \newReading Knowing about these biases helps keep you on guard and better able to seek out the truth. Finally, if you ever find yourself really believing that the world is actually getting worse, remember these thoughts from evolutionary psychologist Steven Pinker:\par
  \noindent\makebox[\textwidth][c]{
    \begin{minipage}[t]{0.9\textwidth}
      Cruelty as entertainment, human sacrifice to indulge superstition, slavery as a labor-saving device, conquest as the mission statement of government, genocide as a means of acquiring real estate, torture and mutilation as routine punishment, the death penalty for misdemeanors and differences of opinion, assassination as the mechanism of political succession, rape as the spoils of war, pogroms as outlets for frustration, homicide as the major form of conflict resolution\textemdash{}all were unexceptional features of life for most of human history. But, today, they are rare to nonexistent in the West, far less common elsewhere than they used to be, concealed when they do occur, and widely condemned when they are brought to light.
    \end{minipage}
  }
  
  With this fourth cup of wine, we understand that bringing freedom to all the world is not some abstract, never ending task\textemdash{}as it was during every past Seder since we left Egypt\textemdash{}but something actually achievable in our lifetimes. Use that as motivation, spread the good news to other people, and work to make universal freedom happen even sooner than predicted.


}